%% SYNTHETIC TEST UNIT generated by writing/tools/make_outline_book.py from FINAL_OUTLINE_v2 -- not book text
\chapter{因子化的破坏与非标准情形}\label{ch:break}
本单元为测试用的合成文本（由 FINAL\_OUTLINE\_v2 自动生成），陈述内容从略。

\section{第 1 节（测试）}\label{sec:break-overview}

\begin{definition}[标准与广义 TMD 因子化]\label{def:break-generalized-tmd}
（测试用合成陈述：内容从略。）
\uses{def:tmd-unsub}
\end{definition}

\begin{proposition}[测量相容性判据 \tagR]\label{prop:break-measurement}
（$R_{\mathrm{cond}}$：在所列编号假设之下证明完整；结论的地位不高于其中最弱的假设。）
前提：\cref{asm:appO-lcopt}、\cref{lem:appO-unitarity}。
（测试用合成陈述：内容从略。）
\uses{lem:appO-unitarity, asm:appO-lcopt}
\end{proposition}

\begin{remark}[失效的位置]\label{rem:break-measurement-failures}
（测试用合成陈述：内容从略。）
\uses{prop:break-measurement}
\end{remark}

\section{第 2 节（测试）}\label{sec:break-tmd}

\begin{proposition}[单个额外胶子的色因子（胶子极强度） \tagR]\label{prop:break-colour}
（$R_{\mathrm{PT}}$：在重整化微扰论中对红外正规化的部分子态严格。）
（测试用合成陈述：内容从略。）
\uses{prop:prelim-wilson-rules, lem:qcd-colour-identities}
\end{proposition}

\begin{example}[由投影算符得到 $C_G$ 表；阿贝尔极限；纠缠检验]\label{ex:break-colour}
（测试用合成陈述：内容从略。）
\uses{prop:break-colour}
\end{example}

\begin{theorem}[强子—强子末态强子的 TMD 因子化不成立 \tagV]\label{thm:break-tmd}
（测试用合成陈述：内容从略。）
\uses{prop:break-colour, def:break-generalized-tmd, rem:break-measurement-failures}
\end{theorem}
\paragraph{地位与缺口}测试：见文献。

\section{第 3 节（测试）}\label{sec:break-spacelike}

\begin{example}[Collins 2007 的双胶子反常（阿贝尔标量模型）]\label{ex:break-c07}
（测试用合成陈述：内容从略。）
\uses{thm:break-tmd}
\end{example}

\begin{proposition}[类空共线极限中严格共线因子化的失效 \tagV]\label{prop:break-spacelike}
（测试用合成陈述：内容从略。）
\uses{prop:sudakov-amplitude-factorization}
\end{proposition}
\paragraph{地位与缺口}测试：见文献。

\begin{example}[CdFR 单圈类空劈裂的色结构]\label{ex:break-spacelike}
（测试用合成陈述：内容从略。）
\uses{prop:break-spacelike}
\end{example}

\section{第 4 节（测试）}\label{sec:break-mpi}

\begin{proposition}[超出 NNLO 的质量奇异性因子化 \tagC]\label{prop:break-mass-singularity}
（测试用合成陈述：内容从略。）
\uses{def:pdf-quark, def:pdf-gluon}
\end{proposition}
\paragraph{地位与缺口}测试：本陈述的地位由\cref{prop:jets-sll}、\cref{prop:break-spacelike}决定。

\begin{proposition}[Coulomb 相位与软交换不对易导致的共线因子化破坏 \tagV]\label{prop:break-fss}
（测试用合成陈述：内容从略。）
\uses{prop:break-spacelike, prop:jets-sll}
\end{proposition}
\paragraph{地位与缺口}测试：本陈述的地位由\cref{prop:jets-sll}决定。

\begin{proposition}[对多重部分子相互作用敏感的观测量 \tagV]\label{prop:break-jet-veto}
（测试用合成陈述：内容从略。）
\uses{prop:break-measurement}
\end{proposition}
\paragraph{地位与缺口}测试：见文献。

\section{第 5 节（测试）}\label{sec:break-diffraction}

\begin{definition}[衍射 DIS 与衍射部分子分布]\label{def:break-diffraction}
（测试用合成陈述：内容从略。）
\uses{def:pdf-quark}
\end{definition}

\begin{theorem}[衍射 DIS 的因子化 \tagA]\label{thm:break-ddis}
（测试用合成陈述：内容从略。）
非微扰地位：\cref{asm:intro-bridge}。
\uses{def:break-diffraction, thm:method-template, prop:method-glauber-deformation, thm:disall-factorization, asm:intro-bridge}
\end{theorem}
\begin{proofsketch}测试。\end{proofsketch}
\paragraph{地位与缺口}测试：本陈述的地位由\cref{thm:disall-factorization}决定。

\begin{proposition}[强子—强子硬衍射中的因子化破坏 \tagV]\label{prop:break-hh-diffraction}
（测试用合成陈述：内容从略。）
\uses{thm:break-ddis}
\end{proposition}
\paragraph{地位与缺口}测试：本陈述的地位由\cref{prop:break-measurement}决定。

\section{第 6 节（测试）}\label{sec:break-dps}

\begin{definition}[双部分子分布与双 Drell–Yan]\label{def:break-dpdf}
（测试用合成陈述：内容从略。）
\uses{def:pdf-quark}
\end{definition}

\begin{theorem}[双 Drell–Yan 的因子化 \tagA]\label{thm:break-double-dy}
（测试用合成陈述：内容从略。）
非微扰地位：\cref{asm:intro-bridge}。
\uses{def:break-dpdf, thm:dy-factorization, asm:intro-bridge}
\end{theorem}
\begin{proofsketch}测试。\end{proofsketch}
\paragraph{地位与缺口}测试：见文献。

\begin{conjecture}[有色末态的双部分子散射 \tagC]\label{conj:break-dps-coloured}
（测试用合成陈述：内容从略。）
\uses{thm:break-double-dy}
\end{conjecture}
\paragraph{地位与缺口}测试：本陈述的地位由\cref{conj:dy-coloured-final}决定。

\section{第 7 节（测试）}\label{sec:break-summary}

\begin{conjecture}["口袋公式" $\sigma_{\rm DPS}=\sigma_A\sigma_B/\sigma_{\rm eff}$ \tagC]\label{conj:break-pocket}
（测试用合成陈述：内容从略。）
\uses{def:break-dpdf}
\end{conjecture}
\paragraph{地位与缺口}测试：本陈述的地位由\cref{thm:break-double-dy}决定。
